\documentclass[10pt,a4paper]{amsart}
\usepackage[margin=19mm]{geometry}
\usepackage{amsmath,amssymb,mathtools}
\usepackage[scr=rsfs,cal=euler]{mathalfa}
\usepackage{xfrac}
\usepackage[unicode,hidelinks,bookmarks=false]{hyperref}
\newcommand{\ie}{i.e.\ }
\newcommand{\cf}{cf.\ }
\newcommand{\resp}{respectively\ }
\newcommand{\Subsection}{\subsection}
\providecommand{\nobreakdash}{\nobreak\hbox{-}}
\setlength{\emergencystretch}{2em}
\multlinegap0pt
\usepackage{amssymb}
\usepackage{mathtools}
\usepackage{tikz}
\usepackage[shortlabels]{enumitem}
\usepackage{algorithm}\usepackage{algpseudocode}
\newtheorem{proposition}{Proposition}
\usepackage{cleveref}
\crefname{proposition}{Proposition}{Propositions}
\title{Continuum dynamics from quantised interaction rules}
\author{Park Junhu, Youngsoo Ha, Myungjoo Kang}
\date{Source: arXiv:2604.06947v2 (CC BY 4.0)}
\begin{document}
\maketitle

\noindent\emph{Source and licence.} Continuum dynamics from quantised interaction rules. Version arXiv:2604.06947v2 (CC BY 4.0), distributed by arXiv under the Creative Commons Attribution 4.0 International licence, http://creativecommons.org/licenses/by/4.0/, as recorded in the arXiv licence metadata for this exact version and shown on the abstract page https://arxiv.org/abs/2604.06947v2. Full licence notice: \texttt{licenses/arxiv-2604.06947-CC-BY-4.0.txt}.\par\smallskip

\noindent\textit{Editorial scope.} Counted unit: the article's proposition prop:operator\_quantisation\_error (source0.tex lines 208--224). The article's complete original proof of it is delivered with it (source0.tex lines 227--241, one \texttt{proof} environment of 15 lines). No mathematics is written, repaired or re-derived: the statement and the proof are the article's own lines, in the article's own notation, and no retained line is substituted or re-flowed.  Retained uncounted as the source-local context the proof needs: lines 142--147; lines 165--168.  Nothing was omitted from any retained span, so the package carries no omission record.  No retained line contains a citation, so the wrapper carries no bibliography and no bibliography entry.  No retained line contains a figure environment, an image-inclusion command or drawing code, so no image is delivered and the package declares no asset dependency.  Editorial formatting only, in addition: an AMS \texttt{amsart} preamble that restates the article's own declarations and macros used by the retained lines, namely the environments \texttt{proposition} on separate counters, together with the standard packages those lines need.  The retained environments print in this document as Proposition 1 in order of appearance, whereas the article prints the same items with the numbering of its own full document.  The complete native source of the article as distributed in the arXiv e-print for this exact version is bundled unchanged as \texttt{sources/198053/source0.tex}.
\begin{equation}
\phi^\pm(q) = \mathrm{round}\!\left(
 f^\pm(\delta q)\frac{\Delta t}{\Delta x}\delta^{-1}
\right).
\label{eq:flux_table}
\end{equation}

\begin{equation}
q_i^{n+1}=q_i^n-\bigl(F_{i+\frac12}^n-F_{i-\frac12}^n\bigr),
\label{eq:update_rule_conservative_form}
\end{equation}

 \begin{proposition}[Update-level correspondence with the classical split scheme]
\label{prop:operator_quantisation_error}
Let $u_i=\delta q_i$. Then the reconstructed one-step FQNM increment differs from the classical split finite-volume increment by $\mathcal O(\delta)$:
\begin{equation}
u_i^{n+1}-u_i^n
= -\frac{\Delta t}{\Delta x}\Bigl(\mathcal F(u_i,u_{i+1})-\mathcal F(u_{i-1},u_i)\Bigr)
+ \mathcal O(\delta).
\label{eq:operator_quantisation_error_bound}
\end{equation}
Equivalently, after division by $\Delta t$,
\begin{equation}
\frac{u_i^{n+1}-u_i^n}{\Delta t}
= \mathcal O_h[u]_i + \mathcal O\!\left(\frac{\delta}{\Delta t}\right).
\label{eq:operator_quantisation_error_asymptotic}
\end{equation}
Under fixed CFL scaling $\Delta t=\mathcal O(\Delta x)$, the operator-level residual is therefore $\mathcal O(\delta/\Delta x)$.
\end{proposition}

\begin{proof}
From \cref{eq:flux_table}, write
\[
\phi^\pm(q)= f^\pm(\delta q)\frac{\Delta t}{\Delta x}\delta^{-1}+\varepsilon^\pm(q),
\qquad |\varepsilon^\pm(q)|\le \tfrac12.
\]
Substituting into \cref{eq:update_rule_conservative_form} and multiplying by $\delta$ yields
\[
u_i^{n+1}-u_i^n
= -\frac{\Delta t}{\Delta x}\Bigl(\mathcal F(u_i,u_{i+1})-\mathcal F(u_{i-1},u_i)\Bigr)
+ \mathcal O(\delta),
\]
since $\delta\varepsilon^\pm(q)=\mathcal O(\delta)$. This gives \cref{eq:operator_quantisation_error_bound}. Dividing by $\Delta t$ gives \cref{eq:operator_quantisation_error_asymptotic}, and under fixed CFL scaling $\Delta t=\mathcal O(\Delta x)$ this becomes $\mathcal O(\delta/\Delta x)$.

\end{proof}

\end{document}
